\chapter{Conclusions}\label{ch:conclusions}

\section{What the model turned out to be}\label{sec:conc1}

The BTW sand pile is usually introduced as a rule about four grains and four neighbours.
The useful object is not the rule but the table it defines. Once the rule is written as
the toppling matrix $\D$ of \eqref{eq:delta}, four separate claims --- that topplings
commute, that an avalanche ends, that its endpoint is unique, and that the number of
recurrent configurations is $\det(-\D)$ --- all follow from four sign conditions on that
table, each constraining one part of the dynamics. The chain runs \eqref{eq:4a}
$\Rightarrow$ commutativity, \eqref{eq:4c} and \eqref{eq:4d} $\Rightarrow$ termination
and uniqueness, and \eqref{eq:4a} and \eqref{eq:4c} $\Rightarrow$ the graph-Laplacian
reading that makes the algebra possible. Nothing in the chain refers to squares.

That reframing is what made the extension cheap. The hexagonal lattice of
Chapter~\ref{ch:hexagonal} needs a new matrix, threshold and distance, but no new proof
of the Abelian property, because the existing proof never uses the lattice. The general
statements in the literature --- the model on an arbitrary directed graph
\cite{dhar1999} and on an arbitrary finite multigraph with the degree as threshold
\cite{jarai2018} --- say the same at the level of generality where the question does not
arise. The corresponding rule for later work is that two things must be checked, and
neither is geometric: the instability test must involve the height of the site alone, and
a toppling must never remove sand from another site. Changes of geometry, degree,
boundary shape or dimension preserve both; changes of \emph{rule} may not. Randomness
alone is not the problem --- the Manna model \cite{manna1990} is still Abelian
\cite{dhar1999} --- which shows that Abelianity and shared exponents are separate
properties.

The stationary state is correspondingly well characterised. The recurrent configurations
carry a uniform measure with weight $1/\absR$, where $\absR = \det(-\D)$ counts the
spanning trees of the wired graph, and that measure fixes the mean height exactly:
$\meanh = 17/8$ on the square lattice, with the four height probabilities known in closed
form. Those are the quantities against which any implementation is checked, because they
involve no fitting, no cut-off and no assumption about a tail.

\section{What the hexagonal comparison settled}\label{sec:conc2}

The comparison was carried out for everything that follows from the two toppling matrices
alone, and it splits cleanly. \emph{Universal}: the exponents $\tau_s = 6/5$ and
$\tau_d = 5/4$, since their derivation routes through the $q \to 0$ Potts model and the
two-dimensional conformal field theory with $c = -2$ \cite{majumdar1992,dhar1999}, which
is sensitive to the dimension and not to the coordination number.
\emph{Lattice-specific and derivable}: the threshold is forced to $6$ by \eqref{eq:4c};
the row-sum classes become $\{0,-2,-3\}$; the mean height becomes $(5+\loopc_{\hx})/2$,
with $17/8$ the case $\loopc = 5/4$; and the sand pile group grows at a different
exponential rate, $h_{\hx} \approx 1.62$ against $4G/\pi = 1.16624$.
\emph{Lattice-specific and affecting the analysis rather than the physics}: the $3/2$
ball-size ratio puts the finite-size cut-off elsewhere, so the collapse variable must be
chosen with care and the two lattices must be compared at matched sizes.

Two features of the change of threshold follow directly from the row sums and are easy to
get wrong. A hexagonal boundary is deeper, not shallower: its six vertices lose three
grains per toppling against two for a square corner, and at equal $N$ the deepest class
carries $30\%$ more of the total loss. Only the relative spread of the per-toppling loss
is narrower. What the hexagonal boundary does have is isotropy --- constant graph
distance from the centre in all six directions, where a square boundary is at $L/2$ along
an axis and $L/\sqrt{2}$ along a diagonal --- and isotropy, not evenness of loss, is the
property that bears on the shape of large avalanches and on the cut-off.

The measurements that would test the universal-exponent prediction are listed in
Table~\ref{tab:hexplan}. The single most informative number is the mean height, since by
\eqref{eq:zetahex} it determines the looping constant of the triangular lattice.
